\documentclass[11pt,a4paper]{article}

% ─── Packages ──────────────────────────────────────────────────────────────────
\usepackage[a4paper, margin=2.5cm]{geometry}
\usepackage{amsmath, amssymb}
\usepackage{graphicx}
\usepackage{booktabs}
\usepackage{array}
\usepackage{caption}
\usepackage{subcaption}
\usepackage{float}
\usepackage{hyperref}
\usepackage{xcolor}
\usepackage{enumitem}
\usepackage{parskip}
\usepackage{microtype}
\usepackage{multirow}
\usepackage[style=ieee, backend=bibtex]{biblatex}
\addbibresource{references.bib}

% ─── Colours ───────────────────────────────────────────────────────────────────
\definecolor{accentblue}{RGB}{30, 100, 180}
\hypersetup{colorlinks=true, linkcolor=accentblue,
            citecolor=accentblue, urlcolor=accentblue}

% ─── Custom commands ───────────────────────────────────────────────────────────
\newcommand{\KL}{\mathrm{KL}}
\newcommand{\JSD}{\mathrm{JSD}}
\newcommand{\todo}[1]{\textcolor{red}{\textbf{[TODO: #1]}}}

% ─── Title block ───────────────────────────────────────────────────────────────
\title{%
  \vspace{-1cm}
  {\large DNN 3rd Year Course Project}\\[0.4em]
  {\LARGE\bfseries Predicting Human Annotator Disagreement\\
  in Image Classification}\\[0.4em]
  {\normalsize Using Soft Labels from CIFAR-10H}
}
\author{%
  \todo{Team Name}\\
  \small\textit{\todo{University / Department}}
}
\date{\today}

% ═══════════════════════════════════════════════════════════════════════════════
\begin{document}
\maketitle
\tableofcontents
\newpage

% ═══════════════════════════════════════════════════════════════════════════════
\section*{Abstract}
\addcontentsline{toc}{section}{Abstract}

Human annotators often disagree when labelling ambiguous images, and this
disagreement carries meaningful information about visual uncertainty.
Standard deep neural networks trained with hard (one-hot) labels discard
this signal entirely.
In this work we frame annotator disagreement prediction as a
\emph{soft-label regression} task: given a CIFAR-10 image, predict the full
probability distribution over classes that a crowd of human annotators would
assign.
We use the CIFAR-10H dataset~\cite{peterson2019human}, which provides
per-image soft-label distributions derived from ${\sim}51$ annotators per
image, and train ResNet-18 backbones adapted for $32{\times}32$ inputs.

We systematically compare three loss functions — KL divergence, Jensen–Shannon
divergence (JSD), and a custom entropy-calibrated KL loss — and evaluate them
on distribution-matching metrics (KL, JSD, cosine similarity) as well as
entropy-prediction quality (Pearson~$r$, Spearman~$r$, Precision@$K$).
Ablation studies examine the effect of backbone initialisation and prediction
head architecture.
Explainability is provided through Grad-CAM visualisations and systematic
failure-case analysis.
Our results show that \todo{brief key finding, e.g.\ ``the custom
entropy-calibrated loss achieves the best entropy-prediction Pearson~$r$ of
\textbf{X.XX} while remaining competitive on distribution-matching''}.

\newpage

% ═══════════════════════════════════════════════════════════════════════════════
\section{Introduction}

Deep learning classifiers are typically trained to predict a single ``correct''
label for every training example.
Yet many real-world images are genuinely ambiguous: a blurry small animal
could be a cat \emph{or} a dog; an aerial vehicle might be confused for a
bird.
When a large crowd of independent annotators labels such images, the resulting
\emph{label distribution} is non-trivial — it encodes collective human
uncertainty about the visual content.

\subsection{Problem Statement}

Given an image $\mathbf{x} \in \mathbb{R}^{32 \times 32 \times 3}$, we wish
to learn a model $f_\theta$ that predicts a probability distribution
$\hat{p} \in \Delta^9$ (the 9-simplex) over the ten CIFAR-10 classes.
The ground-truth target $p$ is the empirical distribution of human annotator
responses, obtained from CIFAR-10H.
Unlike standard classification, the model must capture not only \emph{which}
class is most likely but also \emph{how uncertain} humans are — i.e., the
entropy $H(p)$ of the distribution.

\subsection{Why Annotator Disagreement Matters}

Peterson et al.~\cite{peterson2019human} showed that training on soft labels
derived from human uncertainty improves out-of-distribution generalisation and
model calibration.
Predicting \emph{where} annotators will disagree is valuable for:
\begin{itemize}[noitemsep]
  \item Active learning — prioritising images that are most confusing to humans;
  \item Data quality auditing — flagging ambiguous samples before labelling;
  \item Uncertainty quantification — informing downstream risk-sensitive decisions.
\end{itemize}

\subsection{Contributions}

\begin{enumerate}[noitemsep]
  \item A systematic comparison of three loss functions for soft-label regression
        on CIFAR-10H, including a novel entropy-calibrated KL loss.
  \item Ablation studies over backbone initialisation (random, ImageNet,
        CIFAR-10 pretrained) and prediction head architecture.
  \item Robustness analysis under annotator subsampling and OOD image corruptions.
  \item Qualitative analysis via Grad-CAM and manual inspection of high-entropy
        failure cases.
\end{enumerate}

% ═══════════════════════════════════════════════════════════════════════════════
\section{Related Work}

\subsection{Soft Labels and Label Smoothing}

Label smoothing~\cite{szegedy2016rethinking} replaces hard one-hot targets
with a smoothed distribution $p' = (1-\epsilon)\,p + \epsilon/K$ to improve
calibration.
Soft labels as used here differ: they are \emph{data-driven}, reflecting real
human uncertainty rather than a uniform prior.

\subsection{Knowledge Distillation}

Hinton et al.~\cite{hinton2015distilling} introduced ``dark knowledge'' —
the soft output distributions of a large teacher model carry richer signal
than hard labels.
Our setting is analogous but the ``teacher'' is the crowd of human annotators
rather than another neural network.

\subsection{CIFAR-10H}

Peterson et al.~\cite{peterson2019human} collected CIFAR-10H by recruiting
workers on Amazon Mechanical Turk to label the 10,000 CIFAR-10 test images.
Each image received ${\sim}51$ independent responses.
The resulting $10{,}000 \times 10$ probability matrix captures fine-grained
human uncertainty at scale.
Training on these soft labels was shown to produce models that generalise
better under distribution shift.

% ═══════════════════════════════════════════════════════════════════════════════
\section{Dataset}

\subsection{CIFAR-10H Description}

CIFAR-10H~\cite{peterson2019human} provides a $10{,}000 \times 10$ matrix of
soft-label probability distributions, one row per image, corresponding to the
10,000 images in the CIFAR-10 \emph{test} set.
Each row sums to 1 and encodes the fraction of annotators who assigned each of
the ten class labels.

\subsection{Data Split}

The 10,000 annotated images were split with a fixed random seed
(\texttt{SEED=42}) as follows:

\begin{table}[H]
\centering
\caption{Dataset split sizes.}
\label{tab:splits}
\begin{tabular}{lrr}
\toprule
\textbf{Split} & \textbf{Images} & \textbf{Fraction} \\
\midrule
Train & 6{,}000 & 60\% \\
Validation & 2{,}000 & 20\% \\
Test & 2{,}000 & 20\% \\
\midrule
Total & 10{,}000 & 100\% \\
\bottomrule
\end{tabular}
\end{table}

Indices were saved to \texttt{data/processed/\{train,val,test\}\_idx.npy} to
ensure reproducibility across all experiments.

\subsection{Exploratory Data Analysis}

\subsubsection{Entropy Distribution}

We compute the Shannon entropy $H(p) = -\sum_{k} p_k \log_2 p_k$ for every
image's soft label distribution.
Figure~\ref{fig:entropy_hist} shows that the vast majority of images have very
low entropy, indicating that annotators strongly agreed on a single label.
A long tail of higher-entropy images represents genuinely ambiguous cases.

\begin{figure}[H]
  \centering
  \includegraphics[width=0.72\textwidth]{figures/entropy_histogram.png}
  \caption{%
    Distribution of Shannon entropy across all 10,000 CIFAR-10H images.
    Most images have entropy near zero (near-unanimous annotator agreement),
    while a small fraction reaches entropies above 1.5 bits, indicating
    substantial disagreement.
  }
  \label{fig:entropy_hist}
\end{figure}

\subsubsection{Per-Class Mean Entropy}

Figure~\ref{fig:per_class_entropy} shows mean annotator entropy broken down by
true CIFAR-10 label.
\emph{Deer} and \emph{cat} exhibit the highest mean disagreement, likely due
to visual similarity to other animal classes (\emph{dog}, \emph{bird}).
\emph{Horse} and \emph{automobile} are the easiest to identify, with
consistently low entropy across annotators.

\begin{figure}[H]
  \centering
  \includegraphics[width=0.72\textwidth]{figures/mean_entropy_per_class.png}
  \caption{%
    Mean Shannon entropy per CIFAR-10 class.
    Animal classes (deer, cat, bird, dog) exhibit higher annotator disagreement
    than vehicle classes (automobile, horse, ship, truck),
    reflecting their greater visual ambiguity.
  }
  \label{fig:per_class_entropy}
\end{figure}

\subsubsection{Human Confusion Matrix}

Figure~\ref{fig:confusion} shows the average annotator distribution per true
class — a soft analogue of a confusion matrix.
Each row is normalised to sum to 1.
The strong diagonal confirms that annotators are generally accurate;
off-diagonal mass concentrates on semantically similar classes
(e.g.\ cat/dog, deer/horse, airplane/bird).

\begin{figure}[H]
  \centering
  \includegraphics[width=0.72\textwidth]{figures/human_confusion_matrix.png}
  \caption{%
    Human confusion matrix: each cell $(i,j)$ is the average probability mass
    that annotators assigned to class $j$ when the true label was class $i$.
    High diagonal values (${\geq}0.90$) confirm strong annotator accuracy;
    non-zero off-diagonal entries reveal consistent confusions between
    visually similar categories.
  }
  \label{fig:confusion}
\end{figure}

\subsubsection{Low- and High-Entropy Examples}

Figure~\ref{fig:entropy_examples} contrasts example images at the extremes of
the entropy distribution.
Low-entropy images contain a single clear, centred object, leaving no
ambiguity for annotators.
High-entropy images are cluttered, low-resolution, or depict objects at class
boundaries.

\begin{figure}[H]
  \centering
  \begin{subfigure}[b]{0.38\textwidth}
    \includegraphics[width=\textwidth]{figures/low_entropy_example.png}
    \caption{Low entropy ($H=0.00$): unambiguous ship.}
  \end{subfigure}
  \hspace{1.5em}
  \begin{subfigure}[b]{0.38\textwidth}
    \includegraphics[width=\textwidth]{figures/high_entropy_example.png}
    \caption{High entropy ($H=2.27$): cluttered scene with multiple possible labels.}
  \end{subfigure}
  \caption{%
    Contrasting examples from the entropy extremes of CIFAR-10H.
    The low-entropy image (left) receives unanimous annotator agreement;
    the high-entropy image (right) elicits substantial disagreement due to
    visual ambiguity and scene clutter.
  }
  \label{fig:entropy_examples}
\end{figure}

% ═══════════════════════════════════════════════════════════════════════════════
\section{Methodology}

\subsection{Data Pipeline}

Raw CIFAR-10 test images were paired with their corresponding CIFAR-10H soft
labels by index alignment (no shuffling required, as both datasets share the
same ordering).
Training images were augmented with random horizontal flips and random crops
($32{\times}32$ with padding 4); no colour jitter or rotation was applied, as
these could alter the visual ambiguity properties that drive annotator
disagreement.
Test images used only centre-crop normalisation.
All pixel values were normalised using CIFAR-10 channel statistics:
$\mu = (0.491, 0.482, 0.447)$, $\sigma = (0.202, 0.199, 0.201)$.

\subsection{Model Architecture}

\subsubsection{Backbone}

We use ResNet-18~\cite{he2016deep} as our backbone, adapted for
$32{\times}32$ CIFAR inputs.
The standard ImageNet ResNet-18 applies a $7{\times}7$ convolution with
stride 2 followed by max-pooling, which would reduce $32{\times}32$ feature
maps to $8{\times}8$ after the stem — collapsing too much spatial detail.
We replace this with a $3{\times}3$ stride-1 convolution and remove the
max-pool layer, keeping the spatial resolution at $32{\times}32$ through the
stem.
Global average pooling then reduces the final $4{\times}4$ feature maps to a
512-dimensional feature vector.

\subsubsection{Prediction Head}

Three head architectures are compared in ablation:
\begin{itemize}[noitemsep]
  \item \textbf{Linear}: a single fully-connected layer ($512 \to 10$).
  \item \textbf{MLP}: two-layer network ($512 \to 256 \to 10$) with ReLU and
        Dropout(0.3).
  \item \textbf{Temperature-scaled}: linear layer with a learnable scalar
        temperature $T$ applied before softmax.
\end{itemize}
All heads are followed by a softmax activation, so the model always outputs a
valid probability distribution.

\subsubsection{Architecture Diagram}

Figure~\ref{fig:arch} illustrates the full forward pass.

\begin{figure}[H]
  \centering
  \includegraphics[width=\textwidth]{figures/resnet_cifar_architecture.pdf}
  \caption{%
    Model architecture.
    A $32{\times}32{\times}3$ image enters the adapted Conv1 block
    ($3{\times}3$, stride 1, no max-pool), passes through the four residual
    stages of ResNet-18, is collapsed by global average pooling to a
    512-dimensional feature vector, and is projected to a 10-dimensional
    probability distribution via the prediction head and softmax.
  }
  \label{fig:arch}
\end{figure}

\subsubsection{Parameter Count}

\begin{table}[H]
\centering
\caption{Parameter count for each model variant.}
\label{tab:params}
\begin{tabular}{lrr}
\toprule
\textbf{Model variant} & \textbf{Total params} & \textbf{Trainable params} \\
\midrule
ResNet-18 + Linear head     & \todo{X,XXX,XXX} & \todo{X,XXX,XXX} \\
ResNet-18 + MLP head        & \todo{X,XXX,XXX} & \todo{X,XXX,XXX} \\
ResNet-18 + Temp.-scaled    & \todo{X,XXX,XXX} & \todo{X,XXX,XXX} \\
\bottomrule
\end{tabular}
\end{table}

\subsection{Loss Functions}

All loss functions operate on the predicted distribution
$\hat{p} = f_\theta(\mathbf{x})$ and the target soft label $p$.

\subsubsection{Loss 1 — KL Divergence}

\begin{equation}
  \mathcal{L}_{\KL}(\hat{p},\, p) \;=\;
  \KL(p \,\|\, \hat{p}) \;=\;
  \sum_{k=1}^{K} p_k \log \frac{p_k}{\hat{p}_k}
  \label{eq:kl}
\end{equation}

KL divergence is asymmetric: it heavily penalises cases where the model
assigns near-zero probability to a class that received non-negligible human
votes.
This makes it a natural fit for our task, as missing a class that annotators
frequently chose is a qualitatively worse error than over-spreading mass.
All distributions are clipped to $\min = 10^{-8}$ to avoid $\log(0)$.

\subsubsection{Loss 2 — Jensen--Shannon Divergence}

\begin{equation}
  \mathcal{L}_{\JSD}(\hat{p},\, p) \;=\;
  \frac{1}{2}\,\KL\!\left(p \,\Big\|\, m\right) \;+\;
  \frac{1}{2}\,\KL\!\left(\hat{p} \,\Big\|\, m\right),
  \quad m = \tfrac{1}{2}(p + \hat{p})
  \label{eq:jsd}
\end{equation}

JSD is symmetric and bounded in $[0, \ln 2]$, which makes it more numerically
stable than KL when either distribution has near-zero entries.
It avoids the infinite-penalty issue that arises in $\KL(p\|\hat{p})$ when
$\hat{p}_k \to 0$ for some class $k$ with $p_k > 0$.

\subsubsection{Loss 3 — Custom: Entropy-Calibrated KL}

\begin{equation}
  \mathcal{L}_{\mathrm{custom}}(\hat{p},\, p)
  \;=\;
  \KL(p \,\|\, \hat{p})
  \;+\;
  \alpha \cdot \bigl[H(\hat{p}) - H(p)\bigr]^2,
  \quad H(q) = -\textstyle\sum_k q_k \log q_k
  \label{eq:custom}
\end{equation}

with $\alpha = 0.5$.
The second term adds an explicit \emph{entropy calibration penalty}: beyond
matching the shape of the distribution, the model is penalised for predicting
the wrong level of annotator uncertainty.
A model that identifies the dominant class correctly but underestimates how
spread-out the human votes were incurs an extra penalty proportional to the
squared entropy error.
This is motivated by the observation that the degree of disagreement (not just
its direction) is the primary quantity of interest for downstream applications.

\subsection{Training Protocol}

To ensure fair comparison, \emph{only the loss function varies} across the
three main experiments.
All other components are fixed as shown in Table~\ref{tab:training}.

\begin{table}[H]
\centering
\caption{Fixed training hyperparameters used in all main experiments.}
\label{tab:training}
\begin{tabular}{ll}
\toprule
\textbf{Component} & \textbf{Value} \\
\midrule
Optimiser              & Adam ($\beta_1=0.9$, $\beta_2=0.999$) \\
Learning rate          & $10^{-3}$ \\
Weight decay           & $10^{-4}$ \\
LR schedule            & CosineAnnealingLR \\
Batch size             & 64 \\
Max epochs             & 150 \\
Early stopping         & Patience 15 (monitored on val loss) \\
Augmentation           & RandomHorizontalFlip + RandomCrop(32, pad=4) \\
Backbone init          & CIFAR-10 hard-label pretrained \\
Random seed            & 42 \\
\bottomrule
\end{tabular}
\end{table}

% ═══════════════════════════════════════════════════════════════════════════════
\section{Experiments}

\subsection{Training Curves}

Figures~\ref{fig:curves_kl}--\ref{fig:curves_custom} show the training and
validation loss curves alongside the validation JSD trajectory for each loss
function.

In all three cases the training loss decreases smoothly, demonstrating that
each objective is optimisable.
The validation loss exhibits moderate fluctuation, characteristic of the
small dataset size (6,000 training images).
The KL-trained model converges fastest and achieves the lowest raw loss
values, but operates on a different scale to JSD and the custom loss.
The custom entropy-calibrated model shows the highest raw loss due to the
additional entropy penalty term, but we evaluate models on a common held-out
metric (JSD and entropy correlations) in Section~\ref{sec:eval}.

\begin{figure}[H]
  \centering
  \includegraphics[width=\textwidth]{figures/training_curves_kl.png}
  \caption{%
    Training and validation KL divergence loss (left) and validation JSD
    (right) for the model trained with $\mathcal{L}_{\KL}$.
    The train/val gap stabilises around epoch 10, after which the model is
    selected by early stopping on validation loss.
  }
  \label{fig:curves_kl}
\end{figure}

\begin{figure}[H]
  \centering
  \includegraphics[width=\textwidth]{figures/training_curves_jsd.png}
  \caption{%
    Training and validation JSD loss (left) and validation JSD (right) for
    the model trained with $\mathcal{L}_{\JSD}$.
    The slower convergence compared to KL (requiring ${\sim}27$ epochs before
    early stopping) reflects the bounded, symmetric nature of JSD.
  }
  \label{fig:curves_jsd}
\end{figure}

\begin{figure}[H]
  \centering
  \includegraphics[width=\textwidth]{figures/training_curves_custom.png}
  \caption{%
    Training and validation loss (left) and validation JSD (right) for the
    model trained with $\mathcal{L}_{\mathrm{custom}}$ (entropy-calibrated KL,
    $\alpha=0.5$).
    The higher absolute loss values relative to KL reflect the additional
    entropy penalty term; convergence is faster (early stopping at epoch
    ${\sim}15$) due to the richer gradient signal.
  }
  \label{fig:curves_custom}
\end{figure}

\subsection{Core Evaluation}
\label{sec:eval}

All models are evaluated on the held-out test split (2,000 images) after
loading the best checkpoint (lowest validation loss).
We report six metrics computed over the full test set.

\subsubsection{Metric Definitions}

\paragraph{KL Divergence} $\KL(p \| \hat{p})$ as defined in
Equation~\eqref{eq:kl}. Lower is better.

\paragraph{Jensen--Shannon Divergence} $\JSD(p, \hat{p})$ as defined in
Equation~\eqref{eq:jsd}, computed via \texttt{scipy.spatial.distance.jensenshannon}
(squared). Lower is better.

\paragraph{Cosine Similarity}
$\cos(p, \hat{p}) = p \cdot \hat{p} / (\|p\| \|\hat{p}\|)$.
Higher is better.

\paragraph{Pearson $r$} Pearson correlation between per-image true entropies
$H(p_i)$ and predicted entropies $H(\hat{p}_i)$ across the test set.
Measures whether the model correctly ranks images by their degree of
disagreement. Higher is better.

\paragraph{Spearman $\rho$} Rank-order correlation of the same entropy pairs.
Robust to outliers. Higher is better.

\paragraph{Precision@$K$}
The fraction of the $K$ most disagreement-prone images by predicted entropy
that are truly in the top-$K$ by annotator entropy.
Evaluated at $K \in \{100, 200, 500\}$.

\subsubsection{Results}

\begin{table}[H]
\centering
\caption{%
  Test-set evaluation metrics for each loss function.
  Best result per column is \textbf{bolded}.
  $\downarrow$ = lower is better; $\uparrow$ = higher is better.
}
\label{tab:results}
\setlength{\tabcolsep}{5pt}
\begin{tabular}{lcccccccc}
\toprule
\textbf{Loss} & \textbf{KL}$\downarrow$ & \textbf{JSD}$\downarrow$ &
\textbf{Cos}$\uparrow$ & \textbf{Pearson}$\uparrow$ &
\textbf{Spearman}$\uparrow$ &
\textbf{P@100}$\uparrow$ & \textbf{P@200}$\uparrow$ & \textbf{P@500}$\uparrow$ \\
\midrule
KL Divergence          & \todo{} & \todo{} & \todo{} & \todo{} & \todo{} & \todo{} & \todo{} & \todo{} \\
Jensen--Shannon        & \todo{} & \todo{} & \todo{} & \todo{} & \todo{} & \todo{} & \todo{} & \todo{} \\
Custom (KL + Entropy)  & \todo{} & \todo{} & \todo{} & \todo{} & \todo{} & \todo{} & \todo{} & \todo{} \\
\bottomrule
\end{tabular}
\end{table}

\subsubsection{Entropy Scatter Plot}

\todo{Insert entropy scatter plot (predicted vs.\ true entropy on test set,
with diagonal reference line and Pearson r annotation) once evaluation is
complete. Save as \texttt{figures/entropy\_scatter.png}.}

% Uncomment when figure is ready:
%\begin{figure}[H]
%  \centering
%  \includegraphics[width=0.55\textwidth]{figures/entropy_scatter.png}
%  \caption{%
%    Predicted vs.\ true Shannon entropy for all 2,000 test images under the
%    best-performing model.
%    The dashed red line indicates perfect prediction.
%    Pearson $r = \todo{X.XXX}$ indicates \todo{strong/moderate/weak}
%    correlation between predicted and actual annotator disagreement.
%  }
%  \label{fig:entropy_scatter}
%\end{figure}

\subsubsection{Qualitative Examples}

\todo{%
  Insert a $3 \times 5$ grid of example test images (5 low-, 5 medium-,
  5 high-entropy) once evaluation is complete, each with true and predicted
  distribution bar charts.
  Save as \texttt{figures/qualitative\_examples.png}.
}

% ─── Ablations ─────────────────────────────────────────────────────────────────
\subsection{Ablation Studies}

\subsubsection{Ablation A — Backbone Initialisation}

We train the same model (ResNet-18, linear head, KL loss) under three
initialisation strategies and report test JSD in Table~\ref{tab:ablation_init}.

\begin{table}[H]
\centering
\caption{%
  Effect of backbone initialisation on test JSD.
  The best result is \textbf{bolded}.
}
\label{tab:ablation_init}
\begin{tabular}{lcc}
\toprule
\textbf{Initialisation} & \textbf{JSD}$\downarrow$ & \textbf{Pearson $r$}$\uparrow$ \\
\midrule
Random                 & \todo{} & \todo{} \\
ImageNet pretrained    & \todo{} & \todo{} \\
CIFAR-10 pretrained    & \todo{} & \todo{} \\
\bottomrule
\end{tabular}
\end{table}

\todo{Fill in numbers and add 1--2 sentences of analysis.}

\subsubsection{Ablation B — Prediction Head Architecture}

\begin{table}[H]
\centering
\caption{Effect of prediction head architecture on test JSD (KL loss, CIFAR-10 init).}
\label{tab:ablation_head}
\begin{tabular}{lcc}
\toprule
\textbf{Head} & \textbf{JSD}$\downarrow$ & \textbf{Pearson $r$}$\uparrow$ \\
\midrule
Linear              & \todo{} & \todo{} \\
MLP (2-layer)       & \todo{} & \todo{} \\
Temperature-scaled  & \todo{} & \todo{} \\
\bottomrule
\end{tabular}
\end{table}

\todo{Fill in numbers and add 1--2 sentences of analysis.}

\subsubsection{Ablation C — Training Data Strategy}

\begin{table}[H]
\centering
\caption{%
  Effect of training data strategy on test JSD.
  Strategy 1: train on soft labels only.
  Strategy 2: pretrain on 50k hard labels, then fine-tune on soft labels.
}
\label{tab:ablation_data}
\begin{tabular}{lcc}
\toprule
\textbf{Strategy} & \textbf{JSD}$\downarrow$ & \textbf{Pearson $r$}$\uparrow$ \\
\midrule
Soft labels only             & \todo{} & \todo{} \\
Hard-label pretrain + finetune & \todo{} & \todo{} \\
\bottomrule
\end{tabular}
\end{table}

\todo{Fill in numbers and analysis.}

% ─── Robustness ────────────────────────────────────────────────────────────────
\subsection{Robustness Checks}

\subsubsection{Check A — Annotator Subsampling}

To simulate scenarios with fewer annotators, we re-compute soft labels by
subsampling $k \in \{5, 10, 20, 35, 51\}$ annotators per image from the raw
CIFAR-10H response matrix and measure the JSD between subsampled and
full-annotator distributions.

\todo{Insert annotator subsampling plot (mean JSD vs.\ $k$ annotators).}

\subsubsection{Check B — OOD Corruptions}

We apply Gaussian noise, Gaussian blur, and contrast reduction at five
severity levels to the test images and measure mean predicted entropy at each
severity level.
A well-calibrated model should predict higher entropy as corruptions increase
(more uncertain inputs).

\todo{Insert OOD corruption plot (severity vs.\ mean predicted entropy, one
line per corruption type).}

% ─── Explainability ────────────────────────────────────────────────────────────
\subsection{Explainability Analysis}

\subsubsection{Grad-CAM Visualisations}

Gradient-weighted Class Activation Maps (Grad-CAM)~\cite{selvaraju2017gradcam}
were computed using the last convolutional layer of the ResNet-18 backbone.

\todo{Insert Grad-CAM panel: 5 low-entropy images (left) and 5 high-entropy
images (right), each with the heatmap overlaid. Discuss whether attention is
focused (low entropy) vs.\ diffuse (high entropy).}

\subsubsection{Failure Case Analysis}

\todo{%
  Present the 5--10 images with highest per-image KL divergence between
  prediction and ground truth.
  For each: show image, true distribution, predicted distribution, and
  a brief hypothesis for why the model failed.
}

\subsubsection{Manual Disagreement Source Analysis}

We manually inspected the top 50 highest-entropy test images and assigned each
to one of five categories of disagreement source.

\begin{table}[H]
\centering
\caption{Manual categorisation of the 50 highest-entropy CIFAR-10H images.}
\label{tab:manual}
\begin{tabular}{lrr}
\toprule
\textbf{Category} & \textbf{Count} & \textbf{Percentage} \\
\midrule
Ambiguous identity    & \todo{} & \todo{}\% \\
Poor image quality    & \todo{} & \todo{}\% \\
Multi-object scene    & \todo{} & \todo{}\% \\
Class boundary case   & \todo{} & \todo{}\% \\
Other                 & \todo{} & \todo{}\% \\
\bottomrule
\end{tabular}
\end{table}

\todo{Insert pie chart of category distribution as \texttt{figures/manual\_analysis\_pie.png}.}

% ═══════════════════════════════════════════════════════════════════════════════
\section{Discussion}

\subsection{What Worked}

\todo{%
  Summarise 2--3 findings that held up clearly in experiments.
  E.g.: ``The custom entropy-calibrated loss consistently achieved better
  Pearson $r$ across all three ablation conditions, confirming that explicitly
  supervising the entropy level provides a useful gradient signal beyond
  distribution shape.''
}

\subsection{What Did Not Work and Why}

\todo{%
  Be honest about negative results.
  E.g.: ``The MLP head did not outperform the linear head, suggesting that the
  feature representations from the pretrained backbone are already
  well-linearly-separable for this task.''
  Or: ``Predicted entropy did not consistently increase with OOD corruption
  severity; we hypothesise that BatchNorm statistics make the backbone
  insensitive to additive Gaussian noise.''
}

\subsection{Limitations}

\begin{itemize}[noitemsep]
  \item Training size is limited to 6,000 images; results may not generalise
        beyond CIFAR-10H.
  \item Soft labels are an empirical estimate of the ``true'' human distribution
        and depend on the specific annotator pool and task interface.
  \item The custom loss introduces a hyperparameter $\alpha$ that was not
        exhaustively tuned; optimal $\alpha$ may differ.
  \item \todo{Add any further limitations identified during experiments.}
\end{itemize}

% ═══════════════════════════════════════════════════════════════════════════════
\section{Conclusion}

We have presented a systematic study of soft-label regression for predicting
human annotator disagreement on CIFAR-10H.
Our adapted ResNet-18 architecture, combined with \todo{best loss name},
achieves \todo{key metric value}, demonstrating that neural networks can learn
meaningful signals about image-level visual ambiguity from crowd-sourced label
distributions.

The entropy-calibrated KL loss provides an explicit mechanism to supervise the
\emph{degree} of disagreement, not just its direction, and \todo{shows
improvement / remains competitive} on entropy-prediction metrics.
Grad-CAM analysis reveals that the model \todo{focuses on / struggles with}
\todo{specific finding from GradCAM}.

Future work could explore richer label-distribution models (e.g.\ Dirichlet
priors), larger backbone architectures, or extension to other ambiguous
datasets beyond CIFAR-10H.

% ═══════════════════════════════════════════════════════════════════════════════
\printbibliography

\newpage
\appendix

\section{Hyperparameter Configuration}
\label{app:config}

\begin{table}[H]
\centering
\caption{Complete hyperparameter configuration (\texttt{config.py}).}
\label{tab:config}
\begin{tabular}{ll}
\toprule
\textbf{Parameter} & \textbf{Value} \\
\midrule
\texttt{SEED}                     & 42 \\
\texttt{BATCH\_SIZE}               & 64 \\
\texttt{LR}                       & $10^{-3}$ \\
\texttt{WEIGHT\_DECAY}             & $10^{-4}$ \\
\texttt{EPOCHS}                   & 150 \\
\texttt{EARLY\_STOPPING\_PATIENCE} & 15 \\
\texttt{CIFAR10H\_TRAIN\_SIZE}     & 6{,}000 \\
\texttt{CIFAR10H\_VAL\_SIZE}       & 2{,}000 \\
\texttt{CIFAR10H\_TEST\_SIZE}      & 2{,}000 \\
\texttt{ENTROPY\_ALPHA}           & 0.5 (custom loss only) \\
\bottomrule
\end{tabular}
\end{table}

\section{Additional Figures}
\label{app:figs}

\todo{Place any supplementary figures here (e.g.\ additional Grad-CAM panels,
extended failure case grids, full ablation grouped bar charts).}

\end{document}
